\section{Identified Gap and Positioning of the Thesis}
\label{sec:gaps}

The preceding sections have described a mature research landscape. Cooperative game theory provides the canonical reference for benefit allocation \cite{shapley1953, moncecchi2020}; a set of scalable approximations has made its computation feasible also for large communities \cite{castro2009, cremers2023, alonsopedrero2024, ferrucci2025}; simplified approaches remain attractive for their transparency towards non-expert users \cite{kulmala2021, fioriti2023}; and a group of fairness indicators, imported from other disciplines, is available to evaluate the resulting allocations (Section~\ref{sec:2.3}). The open problem is therefore no longer how to \emph{compute} a fair allocation, but how to \emph{choose} among allocations that can all be computed.

\subsection*{The research need from which this work originates}

The most recent axiomatic analysis evaluates eight allocation methods against seven desirable properties and concludes that \textit{``no computationally efficient CAM satisfies all desirable properties simultaneously, highlighting structural tensions between fairness axioms, incentive compatibility, and ease of implementation''} \cite{gonzalezAsenjo2025}. For one family of methods the tension is a theorem: \textit{``this trade-off is inherent to the entire price-based family''} \cite{gonzalezAsenjo2025}. The authors also state where the result stops. Their model is \textit{``intentionally minimal''}, with \textit{``known and deterministic consumption and generation within each time slot''}, and among the future developments they list the extension of the framework to settings with \textit{``household flexibility, or multi-period coupling''}, which \textit{``would help evaluate the robustness of CAM properties under more realistic conditions''} \cite{gonzalezAsenjo2025}. Of the four developments they list, this thesis takes up the third and, for the use of realistic profiles, the fourth; the first two, on the formal characterisation of the nucleolus and on efficient algorithms for it, are not pursued. A data-driven comparison by the same group points in the same direction: \textit{``no single cost allocation method dominates across all dimensions and configurations''}, and \textit{``formal property compliance alone is insufficient to assess real-world performance and must be complemented by scenario-specific analysis''} \cite{gonzalezAsenjo2026}.

Two questions follow that these statements do not settle. The incompatibilities are proved for price-based methods and for a single time slot; whether comparable tensions arise for the simplified, performance-based and game-theoretic mechanisms, when time is explicit and members differ in the timing of their injections and withdrawals, has not to our knowledge been quantified. And axioms are satisfied or violated, whereas the choice of a rule requires knowing \emph{by how much} it gives up one property to obtain another.

Answering the second question requires a common yardstick, and the literature repeatedly indicates that it is missing. Dynge and Cali observe that the fairness indicators in use were \textit{``adapted for use in the context of local electricity markets''} from other disciplines and that, before their study, \textit{``no research has been found addressing or comparing the performance of these indicators''}; they conclude that \textit{``additional investigation is needed to refine fairness indicators, ensuring their effectiveness across diverse market models''} \cite{dynge2025}. Gasca et al.\ warn that \textit{``fairness metrics derived from other fields are being applied indiscriminately to ECs' context''} \cite{gasca2024}, Volpato et al.\ indicate among future developments the need to \textit{``identify parameters that can quantitatively compare the different benefit distributions''} \cite{volpato2024}, and for Cremers et al.\ \textit{``there are many other fairness concepts that could be explored in energy applications, and it would be relevant to compare their outcome to Shapley values''} \cite{cremers2023}. A recent review reports the lack of \textit{``conceptual alignment between performance metrics''}, used together \textit{``without a systematic articulation of how these criteria interact or conflict''}, and concludes that efficiency, fairness, scalability and regulatory embedding \textit{``are not orthogonal dimensions. Improving one dimension often weakens another''} \cite{chatel2026}. The vocabulary itself is not shared: equity is \textit{``the equal allocation of costs/benefits for all users, regardless of their contribution''} for Gasca et al.\ \cite{gasca2024}, whereas for Dynge and Cali \textit{``Equity commonly refers to rewarding individuals proportionally to individual contribution''} \cite{dynge2025}.

That the choice of rule matters even when the value to be distributed is fixed ex post is argued, in a stylised model with two agents and no hourly profiles, by Awerkin et al., for whom the sharing rule \textit{``is not a mere redistribution device but a determinant of the equilibrium investment pattern''} and the allocation of the collective benefit is \textit{``a regulatory issue in its own right''} \cite{awerkin2026}. This thesis does not study how members respond to the rule; it measures what each rule rewards.

The comparisons that exist are circumscribed. The closest is that of Couraud et al., who evaluate four local-market mechanisms (pro-rata, two forms of glass-filling and a double auction) on four notions of fairness over 250 simulated communities of twenty members, and conclude that \textit{``there is not one distribution approach that addresses all types of fairness''} \cite{couraud2025}. Their result is read from charts, concerns mechanisms operating at a uniform local price, varies only the photovoltaic penetration, and includes neither cooperative-game rules, nor coalitional stability, nor incentive strength; need-based fairness is left out for lack of data. Under regulated incentives for shared energy, studies that confront more than one mechanism rely on a single case study: three members in \cite{volpato2023}, five in \cite{trevisan2024}, ten in \cite{fioriti2023} and a single municipality in \cite{campagna2024}; the study that varies size, prosumer share and user diversity over 64 case studies applies one mechanism only \cite{ferrucci2025}. Gasca et al.\ cross ten approaches with four indices, and conclude that \textit{``the findings are sensitive to user assets and profile changes''} but also that \textit{``the same conclusions apply to other types of ECs with diverse user types''} \cite{gasca2024}: a claim of generality that their design, a single community of seven households observed over one month, does not test.

In our reading, the gap is therefore the following. The test of whether the verdict on the mechanisms depends on the notion of fairness has been run for four local-market mechanisms and four notions \cite{couraud2025}. What the works reviewed do not show is how mechanisms from different families rank on the same data under a common set of indicators that includes, besides equality and proportionality, \emph{coalitional stability} and \emph{incentive strength}; how much is lost on one notion to gain on another, as a quantity and not as a chart; and how the ordering changes when size, prosumer penetration and composition vary systematically, when the value to be distributed is a regulated transfer on the hourly shared energy.

\subsection*{Contribution of the thesis and its extensions}

The thesis addresses this need with a techno-economic simulation model with hourly resolution, in which the mechanisms are evaluated under identical operating conditions and with the same indicators. Starting from this core, the work extends the analysis in three directions, each motivated by the literature.

The first direction concerns the set of mechanisms. The data-driven comparison covers three of the six families of Section~\ref{sec:2.2} and leaves the framework open to \textit{``methods not considered in this paper''} \cite{gonzalezAsenjo2026}, and Gasca et al.\ list among future developments \textit{``additional types of comparisons''} \cite{gasca2024}. Since the proved tensions concern price-based methods, the open question concerns precisely the methods that do not use prices, among them the performance-based ones, designed to reward the temporal alignment between injections and withdrawals \cite{gianaroli2024, giudice2026}. The family based on internal prices and dynamic tariffs (Section~\ref{sec:2.2.4}) is the exception: it presupposes an internal price, which the Italian framework does not provide for, and it is therefore not represented.

The second direction concerns the evaluation layer. Property compliance is not sufficient \cite{gonzalezAsenjo2026}, and the instrument for the complementary analysis is not well established. The work therefore employs a set of indicators rather than a single one, treats their agreement as a result to be observed, and reports separately the indicators defined as a distance from a reference allocation, whose value is by construction optimal for the mechanism that provides the reference. The quantities computed are analogous, not identical, to the axiomatic properties: the work quantifies them, and does not verify the theorem.

The third direction concerns the configurations. The need to vary the composition recurs among the future developments: communities with \textit{``a higher number of members (e.g., tens or more)''} \cite{volpato2024}, \textit{``the different EC compositions (number of members in the UT groups)''} \cite{minutoLanzini2022}, \textit{``the critical set-up''} of a community at which redistribution mechanisms remain applicable \cite{fina2024}, and \textit{``the number and composition of community members''} \cite{giudice2026}. The comparison is therefore repeated over configurations that vary the number of members, the prosumer penetration and the composition by member category.
% TODO: the composition by category is planned but not yet run (run time, README Section 12): add the configurations before submission or delete the last clause.

\subsection*{Research question and empirical context}

The gap and its extensions converge into a single research question:

\begin{quote}
\emph{When a year of hourly profiles is allocated by mechanisms from different families, (i) does the ranking of the mechanisms change with the notion of fairness adopted (equality, proportionality to the physical contribution, coalitional stability, inclusion) and with the size, the prosumer penetration and the composition by category of the community; (ii) how large are the trade-offs among these notions, including the incentive strength of the rule, along the Pareto frontier; and (iii) which minimal subset of indicators suffices to discriminate among the mechanisms?}
\end{quote}

Each term corresponds to a quantity computed by the model. Equity is measured by indicators of uniformity, of proportionality to the physical contribution, of inclusion and of coalitional stability, all evaluated on the same allocations; stability, available up to eighteen members, is the largest excess of the blocking coalitions over the value to be distributed. The incentive strength of a rule is measured ex post, as the cross-sectional covariance between the share received by each member and the alignment of that member's profile with the community, normalised on a reference mechanism; it indicates how much a rule rewards synchrony, not how members would respond to it, a behavioural response that lies outside the scope of the model. The agreement among indicators is the rank correlation (Kendall's $\tau_b$) between the rankings they induce, the dependence on the configuration is the frequency with which the ordering of a pair of mechanisms reverses along the size and penetration axes, and the minimal subset is obtained by grouping the indicators according to their rank distance and retaining one representative per group. The economic sustainability of participation, computed as the net present value of each member's investment where the investment shares are specified, is treated as a complementary analysis.

The analysis is carried out within the Italian regulatory framework, which is the empirical context of the work and not its motivation. It is an operative and documented instance of a regime in which the value to be distributed is a regulated transfer, not an arbitrage: the shared quantity is the hourly minimum between aggregate injections and withdrawals and the unit remuneration is common to all members, so that the price-dispersion channel that the data-driven comparison identifies as a driver of its results is absent, and members differ only in the temporal alignment of their profiles and, where the regulation provides for it, in their legal category \cite{MASE:Decreto:CER:2023, GSE:CACER:RegoleOperative:2025}. This is the simplest setting in which to pose the question. The admissibility of each allocation is verified ex post against the threshold of the framework and reported.
% TODO: verify this paragraph against the decree text and the GSE Operating Rules (threshold on the premium, reduction factor, exemption of natural persons); on the existing configurations the threshold does not bind (README Section 7.4), so either check it on all 28 configurations or reduce the last sentence to a note.

The thesis answers the question through a simulation model organised in two layers. The valorisation layer has no degrees of freedom: starting from the hourly profiles of the members, it determines the value to be distributed and the constraints acting on it. The allocation layer is the object of the study, and implements sixteen mechanisms belonging to five of the six families of Section~\ref{sec:2.2}, all distributing the same total value and evaluated with the same indicators. The expected outcome is not the identification of an optimal mechanism, but the structure of the interaction among equity, stability and incentive strength: the configurations in which the ranking is reversed, and those in which the indicators agree or diverge. Because several indicators are maximised by construction by a specific mechanism, the quantity of interest is the loss incurred by the other mechanisms on the remaining indicators, rather than the identity of the winner on each indicator.
% TODO: check the number of mechanisms against the final list (the sixteenth runs on placeholder data and must not produce results).
